% GENERATED FILE. Edit canonical lessons, then run npm run generate:books.
% canonical-source-hash: fnv1a64:be9c7a9a04443fa6
% canonical-lessons: ML-C546-kulippani, ML-C546-upajivanam, ML-C546-jolibharam, ML-C546-pravrttiparicayam, ML-C546-nayikkuka

\chapter{Industry and Stress}
\label{ch:ml-industry-and-stress}
\hlchaptermodality{546}

\begin{chapteropening}
\textbf{By the end of this chapter:} \emph{I can say manual labour, a livelihood, a workload, work experience and to lead.}

Complete the last lesson of chapter 546: I can say manual labour, a livelihood, a workload, work experience and to lead.
\end{chapteropening}

\section[\texorpdfstring{\ml{കൂലിപ്പണി} (kūlippaṇi)}{കൂലിപ്പണി (kūlippaṇi)}]{\ml{കൂലിപ്പണി} (kūlippaṇi) — manual labour}

\label{lesson:ML-C546-kulippani}



\begin{quote}
Before the new one: say the Malayalam for planning, then the Malayalam for to cooperate.
\end{quote}

\subsection*{You'll want to know: \ml{കൂലിപ്പണി}}
\textbf{\ml{കൂലിപ്പണി}} — \emph{kūlippaṇi} — \textquotedblleft{}manual labour\textquotedblright{}.

\begin{quote}
\ml{ജോലി} \ml{കിട്ടുന്നതുവരെ} \ml{അവൻ} \ml{കൂലിപ്പണി} \ml{ചെയ്തു}.
\end{quote}

\begin{quote}
\emph{jōli kiṭṭunnatuvare avan kūlippaṇi ceytu.}
\end{quote}

\begin{quote}
Until he found a job, he did daily-wage labour.
\end{quote}

\ml{കൂലി} a daily wage + \ml{പണി} work.

\begin{cousinweb}
\textbf{\ml{കൂലിപ്പണി}} (\emph{kūlippaṇi}) is native Dravidian, a compound of \ml{കൂലി}, \textquotedblleft{}hire, daily wages\textquotedblright{} (DEDR 1905), and \ml{പണി}, \textquotedblleft{}work\textquotedblright{} (DEDR 3884).
\end{cousinweb}

\begin{grammarlens}[title={What you've built}]
One more word for describing what you lived and felt.
\end{grammarlens}

\subsection*{Your turn}
\begin{itemize}
\raggedright
  \item \emph{Say it:} \emph{kūlippaṇi}
  \item \emph{Say it:} \emph{kūlippaṇi}, once more
  \item \emph{Say it:} \emph{jōli kiṭṭunnatuvare avan kūlippaṇi ceytu.} — Until he found a job, he did daily-wage labour.
  \item \emph{Recall:} say the Malayalam for planning, then the Malayalam for to cooperate, then say \emph{kūlippaṇi} again
\end{itemize}

\begin{tcolorbox}[breakable,colback=teal!4,colframe=teal!35!black,title={Before you move on}]
What does \ml{കൂലിപ്പണി} mean? (\textquotedblleft{}Manual labour\textquotedblright{}.) Say it once more.
\end{tcolorbox}
\section[\texorpdfstring{\ml{ഉപജീവനം} (upajīvanaṁ)}{ഉപജീവനം (upajīvanaṁ)}]{\ml{ഉപജീവനം} (upajīvanaṁ) — a livelihood}

\label{lesson:ML-C546-upajivanam}



\begin{quote}
Before the new one: say the Malayalam for to cooperate, then the Malayalam for manual labour.
\end{quote}

\subsection*{You'll want to know: \ml{ഉപജീവനം}}
\textbf{\ml{ഉപജീവനം}} — \emph{upajīvanaṁ} — \textquotedblleft{}a livelihood\textquotedblright{}.

\begin{quote}
\ml{മീൻപിടിത്തമാണ്} \ml{അവരുടെ} \ml{ഉപജീവനം}.
\end{quote}

\begin{quote}
\emph{mīnpiṭittamāṇŭ avaruṭe upajīvanaṁ.}
\end{quote}

\begin{quote}
Fishing is their livelihood.
\end{quote}

\ml{ഉപജീവനമാർഗ്ഗം}, a means of living.

\begin{cousinweb}
\textbf{\ml{ഉപജീവനം}} (\emph{upajīvanaṁ}) is borrowed from Sanskrit \emph{upajīvana}, \textquotedblleft{}livelihood, subsistence\textquotedblright{} (Monier-Williams), from \emph{upa-jīv}, \textquotedblleft{}to live on\textquotedblright{}.
\end{cousinweb}

\begin{grammarlens}[title={What you've built}]
One more word for describing what you lived and felt.
\end{grammarlens}

\subsection*{Your turn}
\begin{itemize}
\raggedright
  \item \emph{Say it:} \emph{upajīvanaṁ}
  \item \emph{Say it:} \emph{upajīvanaṁ}, once more
  \item \emph{Say it:} \emph{mīnpiṭittamāṇŭ avaruṭe upajīvanaṁ.} — Fishing is their livelihood.
  \item \emph{Recall:} say the Malayalam for to cooperate, then the Malayalam for manual labour, then say \emph{upajīvanaṁ} again
\end{itemize}

\begin{tcolorbox}[breakable,colback=teal!4,colframe=teal!35!black,title={Before you move on}]
What does \ml{ഉപജീവനം} mean? (\textquotedblleft{}A livelihood\textquotedblright{}.) Say it once more.
\end{tcolorbox}
\section[\texorpdfstring{\ml{ജോലിഭാരം} (jōlibhāraṁ)}{ജോലിഭാരം (jōlibhāraṁ)}]{\ml{ജോലിഭാരം} (jōlibhāraṁ) — a workload}

\label{lesson:ML-C546-jolibharam}



\begin{quote}
Before the new one: say the Malayalam for manual labour, then the Malayalam for a livelihood.
\end{quote}

\subsection*{You'll want to know: \ml{ജോലിഭാരം}}
\textbf{\ml{ജോലിഭാരം}} — \emph{jōlibhāraṁ} — \textquotedblleft{}a workload\textquotedblright{}.

\begin{quote}
\ml{മാസാവസാനം} \ml{ഓഫീസിൽ} \ml{ജോലിഭാരം} \ml{കൂടും}.
\end{quote}

\begin{quote}
\emph{māsāvasānaṁ ōphīsil jōlibhāraṁ kūṭuṁ.}
\end{quote}

\begin{quote}
At the end of the month the workload at the office goes up.
\end{quote}

\ml{ജോലി} work + \ml{ഭാരം} weight; \ml{ഭാരമുള്ള}, heavy, is already familiar.

\begin{cousinweb}
\textbf{\ml{ജോലിഭാരം}} (\emph{jōlibhāraṁ}) is a Malayalam compound built from \ml{ജോലി}, \textquotedblleft{}work\textquotedblright{}, and \ml{ഭാരം}, from Sanskrit \emph{bhāra}, \textquotedblleft{}a burden, a weight\textquotedblright{} (CDIAL 9459).
\end{cousinweb}

\begin{grammarlens}[title={What you've built}]
One more word for describing what you lived and felt.
\end{grammarlens}

\subsection*{Your turn}
\begin{itemize}
\raggedright
  \item \emph{Say it:} \emph{jōlibhāraṁ}
  \item \emph{Say it:} \emph{jōlibhāraṁ}, once more
  \item \emph{Say it:} \emph{māsāvasānaṁ ōphīsil jōlibhāraṁ kūṭuṁ.} — At the end of the month the workload at the office goes up.
  \item \emph{Recall:} say the Malayalam for manual labour, then the Malayalam for a livelihood, then say \emph{jōlibhāraṁ} again
\end{itemize}

\begin{tcolorbox}[breakable,colback=teal!4,colframe=teal!35!black,title={Before you move on}]
What does \ml{ജോലിഭാരം} mean? (\textquotedblleft{}A workload\textquotedblright{}.) Say it once more.
\end{tcolorbox}
\section[\texorpdfstring{\ml{പ്രവൃത്തിപരിചയം} (pravṛttiparicayaṁ)}{പ്രവൃത്തിപരിചയം (pravṛttiparicayaṁ)}]{\ml{പ്രവൃത്തിപരിചയം} (pravṛttiparicayaṁ) — work experience}

\label{lesson:ML-C546-pravrttiparicayam}



\begin{quote}
Before the new one: say the Malayalam for a livelihood, then the Malayalam for a workload.
\end{quote}

\subsection*{You'll want to know: \ml{പ്രവൃത്തിപരിചയം}}
\textbf{\ml{പ്രവൃത്തിപരിചയം}} — \emph{pravṛttiparicayaṁ} — \textquotedblleft{}work experience\textquotedblright{}.

\begin{quote}
\ml{ആ} \ml{ജോലിക്ക്} \ml{മൂന്ന്} \ml{വർഷത്തെ} \ml{പ്രവൃത്തിപരിചയം} \ml{വേണം}.
\end{quote}

\begin{quote}
\emph{ā jōlikkŭ mūnnŭ varṣatte pravṛttiparicayaṁ vēṇaṁ.}
\end{quote}

\begin{quote}
That job needs three years of work experience.
\end{quote}

\ml{പ്രവൃത്തി} work + \ml{പരിചയം} experience, already familiar.

\begin{cousinweb}
\textbf{\ml{പ്രവൃത്തിപരിചയം}} (\emph{pravṛttiparicayaṁ}) is a learned compound built from Sanskrit \emph{pravṛtti}, \textquotedblleft{}moving onwards, activity\textquotedblright{} (Monier-Williams), and \emph{paricaya}, \textquotedblleft{}acquaintance, practice\textquotedblright{} (CDIAL 7807).
\end{cousinweb}

\begin{grammarlens}[title={What you've built}]
One more word for describing what you lived and felt.
\end{grammarlens}

\subsection*{Your turn}
\begin{itemize}
\raggedright
  \item \emph{Say it:} \emph{pravṛttiparicayaṁ}
  \item \emph{Say it:} \emph{pravṛttiparicayaṁ}, once more
  \item \emph{Say it:} \emph{ā jōlikkŭ mūnnŭ varṣatte pravṛttiparicayaṁ vēṇaṁ.} — That job needs three years of work experience.
  \item \emph{Recall:} say the Malayalam for a livelihood, then the Malayalam for a workload, then say \emph{pravṛttiparicayaṁ} again
\end{itemize}

\begin{tcolorbox}[breakable,colback=teal!4,colframe=teal!35!black,title={Before you move on}]
What does \ml{പ്രവൃത്തിപരിചയം} mean? (\textquotedblleft{}Work experience\textquotedblright{}.) Say it once more.
\end{tcolorbox}
\section[\texorpdfstring{\ml{നയിക്കുക} (nayikkuka)}{നയിക്കുക (nayikkuka)}]{\ml{നയിക്കുക} (nayikkuka) — to lead; to guide}

\label{lesson:ML-C546-nayikkuka}



\begin{quote}
Before the new one: say the Malayalam for a workload, then the Malayalam for work experience.
\end{quote}

\subsection*{You'll want to know: \ml{നയിക്കുക}}
\textbf{\ml{നയിക്കുക}} — \emph{nayikkuka} — \textquotedblleft{}to lead; to guide\textquotedblright{}.

\begin{quote}
\ml{പരിചയമുള്ള} \ml{ആ} \ml{മാനേജർ} \ml{ടീമിനെ} \ml{നന്നായി} \ml{നയിച്ചു}.
\end{quote}

\begin{quote}
\emph{paricayamuḷḷa ā mānējar ṭīmine nannāyi nayiccu.}
\end{quote}

\begin{quote}
That experienced manager led the team well.
\end{quote}

Past: \ml{നയിച്ചു}, \emph{nayiccu}.

A leader is \ml{നേതാവ്}, from the same Sanskrit root.

\begin{cousinweb}
\textbf{\ml{നയിക്കുക}} (\emph{nayikkuka}) is built from Sanskrit \emph{naya}, from \emph{nayati}, \textquotedblleft{}leads\textquotedblright{} (CDIAL 6966), with the Malayalam verb ending -\ml{ിക്കുക}.
\end{cousinweb}

\begin{grammarlens}[title={What you've built}]
One more word for describing what you lived and felt.
\end{grammarlens}

\subsection*{Your turn}
\begin{itemize}
\raggedright
  \item \emph{Say it:} \emph{nayikkuka}
  \item \emph{Say it:} \emph{nayikkuka}, once more
  \item \emph{Say it:} \emph{paricayamuḷḷa ā mānējar ṭīmine nannāyi nayiccu.} — That experienced manager led the team well.
  \item \emph{Recall:} say the Malayalam for a workload, then the Malayalam for work experience, then say \emph{nayikkuka} again
\end{itemize}

\begin{tcolorbox}[breakable,colback=teal!4,colframe=teal!35!black,title={Before you move on}]
What does \ml{നയിക്കുക} mean? (\textquotedblleft{}To lead; to guide\textquotedblright{}.) Say it once more.
\end{tcolorbox}
